\documentclass{article}

\usepackage[utf8]{inputenc}
\usepackage{amsmath,amssymb}
\usepackage{graphicx}
\usepackage{hyperref}

\title{Cartograph: schizoanalytic diagnosis of foundation models}
\author{hinanohart}
\date{2026 (draft)}

\begin{document}
\maketitle

\begin{abstract}
We operationalise Félix Guattari's four schizoanalytic functors — Flows ($F$),
Machinic Phyla ($\Phi$), Existential Territories ($T$), Universes of Reference
($U$) — as diagnostic measurements on foundation models. Phase~1a ships $\Phi$
(SAE co-activation graph) and $U$ (Pareto frontier over multi-objective loss);
Phase~1b extends to $F$ (persistent homology of the velocity field) and $T$
(attention-rollout territories with an SSM control). The work is not a
re-branding of mechanistic interpretability: each functor is gated by a
construct-validity peer review that asks whether the metric \emph{measures} the
Guattarian concept it claims to operationalise.
\end{abstract}

\section{Introduction}
Draft.

\section{Background: schizoanalytic cartographies}
Draft. Cite Guattari (1989), Segovia (2025), Crowe (2019).

\section{Method}
\subsection{$\Phi$ — Machinic Phyla}
\subsection{$U$ — Universes of Reference}
\subsection{$F$ — Flows (Phase 1b)}
\subsection{$T$ — Existential Territories (Phase 1b)}

\section{Experiments}
Draft. Refer to figures in \texttt{paper/figures/}.

\section{Limitations}
Construct validity; PH scale (W1); SAE polysemanticity (W2); T rebrand risk (W3).

\section{Related work}
Pasquinelli; Parisi; Hui; Sha Xin Wei; Crowe; Segovia; Bricken et al.;
Carriere et al.

\bibliographystyle{plain}
% \bibliography{refs}

\end{document}
